\documentclass[10pt,a4paper]{amsart}
\usepackage[margin=19mm]{geometry}
\usepackage{amsmath,amssymb,mathtools}
\usepackage[scr=rsfs,cal=euler]{mathalfa}
\usepackage{xfrac}
\usepackage[unicode,hidelinks,bookmarks=false]{hyperref}
\newcommand{\ie}{i.e.\ }
\newcommand{\cf}{cf.\ }
\newcommand{\resp}{respectively\ }
\newcommand{\Subsection}{\subsection}
\providecommand{\nobreakdash}{\nobreak\hbox{-}}
\setlength{\emergencystretch}{2em}
\multlinegap0pt
\usepackage{amssymb}
\newtheorem{proposition}{Proposition}
\def\Z {{\mathbb{Z}}}
\def\aut{{\mathrm{Aut}}}
\def\lh{{\mathrm{lh}}}
\renewcommand{\restriction}{\mathord{\upharpoonright}}
\title{On the complexity of extensions of non-archimedean Polish groups admitting a compatible complete left-invariant metric}
\author{Longyun Ding, Xu Wang}
\date{Source: arXiv:2605.24379v1 (CC BY 4.0)}
\begin{document}
\maketitle

\noindent\emph{Source and licence.} On the complexity of extensions of non-archimedean Polish groups admitting a compatible complete left-invariant metric. Version arXiv:2605.24379v1 (CC BY 4.0), distributed by arXiv under the Creative Commons Attribution 4.0 International licence, http://creativecommons.org/licenses/by/4.0/, as recorded in the arXiv licence metadata for this exact version and shown on the abstract page https://arxiv.org/abs/2605.24379v1. Full licence notice: \texttt{licenses/arxiv-2605.24379-CC-BY-4.0.txt}.\par\smallskip

\noindent\textit{Editorial scope.} Counted unit: the article's proposition autZw (source0.tex lines 918--924). The article's complete original proof of it is delivered with it (source0.tex lines 926--981, one \texttt{proof} environment of 56 lines). No mathematics is written, repaired or re-derived: the statement and the proof are the article's own lines, in the article's own notation, and no retained line is substituted or re-flowed.  No further source-local context is needed: every cross-reference of every retained line resolves inside the two delivered spans.  Nothing was omitted from any retained span, so the package carries no omission record.  No retained line contains a citation, so the wrapper carries no bibliography and no bibliography entry.  No retained line contains a figure environment, an image-inclusion command or drawing code, so no image is delivered and the package declares no asset dependency.  Editorial formatting only, in addition: an AMS \texttt{amsart} preamble that restates the article's own declarations and macros used by the retained lines, namely the environments \texttt{proposition} on separate counters, together with the standard packages those lines need.  The retained environments print in this document as Proposition 1 in order of appearance, whereas the article prints the same items with the numbering of its own full document.  The complete native source of the article as distributed in the arXiv e-print for this exact version is bundled unchanged as \texttt{sources/201271/source0.tex}.
\begin{proposition}\label{property of autZw}
    \begin{enumerate}
        \item $(L_n)$ is an open neighborhood basis of $1_L$;
        \item $(\aut(\Z^\omega),\mathcal{T})$ is a topological group;
        \item the map $(f,x)\mapsto f(x)$ from $\aut(\Z^\omega)\times\Z^\omega$ to $\Z^\omega$ is continuous.
    \end{enumerate}
\end{proposition}

\begin{proof}
    (1) We first prove the following claim:

{\sl Claim.} For any compact $K\subseteq\Z^\omega$ and open $V\subseteq\Z^\omega$ with $K\subseteq V$, there exists an $m\in\omega$ such that $L_m\subseteq C(K,V)$.

   \begin{proof}[Proof of Claim]
        Since $K$ is compact, we can find $s_0,\dots,s_k\in\Z^{<\omega}$ such that $K\subseteq \bigcup_{i\leq k}N_{s_i}\subseteq V$. Let $m=\max\{\lh(s_i):i\leq k\}$. Then for any $f\in L_m$ and $x\in K$, we can see that
        \begin{align*}
            f(x)\restriction m=\left(\begin{pmatrix}
                I_m & O\\
                \multicolumn{2}{c}{*}
            \end{pmatrix}\begin{pmatrix}
                x\restriction m\\
                *
            \end{pmatrix}\right)\restriction m=x\restriction m.
        \end{align*}
        From $x\in K$ we have $x\in\bigcup_{i\leq k}N_{s_i}$. So $N_{x\restriction m}\subseteq\bigcup_{i\leq k}N_{s_i}$. This implies $f(x)\in\bigcup_{i\leq k}N_{s_i}\subseteq V$. Therefore, we obtain that $L_m\subseteq C(K,V)$.
    \end{proof}

    Now, for any open neighborhood $C(K_0',V_0)\cap\cdots\cap C(K_n',V_n)$ of $1_L$, by the claim above we can find an $m\in\omega$ such that
    $$L_m\subseteq C(K_0',V_0)\cap\cdots\cap C(K_n',V_n).$$
    Then by definition of $L_m$ we can see that $(L_n)$ is an open neighborhood basis of $1_L$.

    (2) For compact $K\subseteq\Z^\omega$, open $V\subseteq\Z^\omega$ and $f\in\aut(\Z^\omega)$, we can see that $C(K,V)f=C(f^{-1}(K),V)$ and $fC(K,V)=C(K,f(V))$. Note that $f^{-1}(K)$ is compact, and $f(V)$ is open. Then a straightforward verification shows that both $(L_nf)$ and $(fL_n)$ are open neighborhood bases of $f$.

    Now, let $f,g\in\aut(\Z^\omega)$ and $n\in\omega$. Suppose $g=(a_{ij})$. Since the matrix $g$ is row-finite, we can find an $M\in\omega$ such that
    $$\forall i<n\,\forall j\geq M\left(a_{ij}=0\right).$$
    Hence we can divide the matrix $g$ as
    \begin{equation*}
        g=\begin{pmatrix}
            A & O\\
            \multicolumn{2}{c}{*}
        \end{pmatrix},
    \end{equation*}
    where $A$ is a matrix of size $n\times M$. Next, we divide the matrix $g^{-1}$ as
    \begin{equation*}
        g^{-1}=\begin{pmatrix}
            B & C\\
            \multicolumn{2}{c}{*}
        \end{pmatrix},
    \end{equation*}
    where $B$ is a matrix of size $M\times n$. Then we have $AB=I_n$ and $AC=O$. By a straightforward computation, we can see that $gL_Mg^{-1}\subseteq L_n$. Thus $gL_M\subseteq L_ng$. Consequently, we obtain that
    $$(L_ng)(fL_M)^{-1}=L_ngL_Mf^{-1}=L_n(gL_Mg^{-1})gf^{-1}\subseteq L_n^2gf^{-1}=L_ngf^{-1}.$$
    Since $n$ is arbitrary, we can see that the map $(g,f)\mapsto gf^{-1}$ is continuous, which means that $(\aut(\Z^\omega),\mathcal{T})$ is a topological group.

    (3) For $n\in\omega$, let
    $$H_n=\prod_{i<n}\{0\}\times\prod_{i\geq n}\Z.$$
    Let $f\in\aut(\Z^\omega)$ and $x\in\Z^\omega$. Suppose $f=(a_{ij})$. Then for any $n\in\omega$, there exists an $M\in\omega$ such that
    $$\forall i< n\,\forall j\geq M\left(a_{ij}=0\right).$$
    Then a direct computation shows that
    \begin{align*}
        (L_nf)(x+H_M)&=\{g(y):g\in L_nf\land y\in x+H_M\}\\
        &\subseteq f(x)+H_n.
    \end{align*}
    Hence the map $(f,x)\mapsto f(x)$ is continuous.
\end{proof}

\end{document}
